\section{Finite cohomology for a class supported on one arc}

% BEGIN H82A1
\paragraph{H82A1. A crystalline comparison in the presence of several cyclotomic poles.}
Let $K/\mathbb Q_p$ be unramified, $p$ odd, with the level-one
data above. On the original Wach variable $X$ put
$q_k=\varphi^k Q$, $Q=\varphi(X)/X$. Suppose an integral Herr
pair $(y,z)$ with values in the \emph{\'etale} module is invariant
under the tame group, $z\in XN$, and some finite product of the
$q_k$ takes $y$ into $N$. This pair represents a finite integral
class. We prove the assertion with its given integral representatives.

Restrict to $K(\mu_p)$ and choose B81L2's triangular Witt basis,
the old Wach columns followed by $e+b$ with
$(\varphi-1)b=y$ in the Witt embedding. Its Galois representation
is the extension from the integral Herr pair by LWR5. As in B81L2
the basis vectors are fixed by the cyclotomic kernel; the generator
preserves their series span and acts trivially modulo $X$ since
$z\in XN$. Its Frobenius matrix $P$ and its inverse have only a
finite product of cyclotomic factors in their denominators, by
triangular inversion and the ordinary height bounds on the old
block. Choose $k_0,d\ge0$ large enough that
$P_0^dP,P_0^dP^{-1}$ are integral, where
\[
       P_0=\prod_{k=0}^{k_0}q_k,\qquad
       D=\prod_{k=0}^{k_0}\varphi^k X,\qquad
       \varphi(D)=DP_0 .
\]
If $Y$ is the matrix of the triangular Witt basis in constant
lattice coordinates, then $\varphi(Y)=YP$. Thus $D^dY$ and
$D^dY^{-1}$ satisfy the integral Frobenius equations of L78W2,
using $P_0^dP$ on the right and $P_0^dP^{-1}$ on the left,
respectively. They belong to the integral Witt ring with
coefficients. The image of $X=[\epsilon^{1/p}]-1$ in
$B_{\rm cris}$ is a unit by L78W4. Frobenius on this ring is
unital, so each $\varphi^k X$ is a unit as well. This evaluates
$Y,Y^{-1}$ as mutually inverse matrices there.

On the series side solve for the normalized generator-invariant
frame $U$ and its inverse near zero by the coefficient recursion
of L78W9. Every $q_k(0)=p$ is nonzero over the coefficient field.
Since Frobenius and the generator commute, the same uniqueness
argument gives $P\varphi(U)=U C$ with constant invertible $C$,
cyclically over the letters. Let $t(X)=\log(1+X)$ and
$D'=\prod_{k=0}^{k_0}\varphi^k(t(X)/X)$. It is analytic on the
open unit disc and
$P_0\varphi(D')=p^{k_0+1}D'$.
Hence $(D')^d U$ satisfies a Frobenius equation with analytic
matrix coefficients on the whole disc, using $P_0^dP$.
The inverse transpose satisfies the analogous equation using
$P_0^dP^{-\mathsf T}$. Iterating Frobenius towards zero as in
L78W10 continues both matrices throughout the disc.
The product $D D'$ is a nonzero rational scalar times
$t(X)^{k_0+1}$, so one logarithmic power clears all their poles.
L78W6--L78W7 and L78W11 give crystalline evaluation and
equivariance on this localization, with $t(X)$ invertible.
The evaluated matrix $YU$
therefore gives a basis of invariant periods of the extension.
This proves finiteness after restriction. Corestriction and
division by $[K(\mu_p):K]=p-1$ preserve the class and its finite
condition, finishing the assertion over $K$.
% END H82A1

% BEGIN H82A2
\paragraph{H82A2. The class of a cut vector and its dependence on the residue.}
Use $v$ for the tame invariant parameter of L81S2, $L_i=N_i^\Delta$,
and put $P_k=\varphi^k(v+p)$. As series in $X$ they are associates
of the $q_k$ of H82A1. Write $L(c,j)$ for the series span of
$\varphi^j L_c$ at letter $c+j$ and
$D_j=v^M\prod_{k<j}P_k^{3g}$. The source's cut space at $c$
consists of vectors $w$ with $(\gamma-1)w\in vL_c$ and
$\Lambda w$ in the span obtained by applying Frobenius from the
preceding cut. Here $\Lambda$ is the prescribed product of
the $P_k$ along that arc. A pair for $w$ in $\mathrm{IsRho}$ has
a first entry supported on $c+j$ with value
$p^\sigma\varphi^j w-\tau$, $\tau\in L(c,j)$. Its second entries
are in $v^M L_i$, and its two differences satisfy the Herr equation.
Successive height inclusions show that multiplying
$L(c,j)$ by a finite product of the $P_k$ puts it in $L_{c+j}$:
at each step, multiply by a sufficient power of $P_0$ after
applying Frobenius to the factors from the preceding step.
The preceding-cut condition gives the same finite-pole property
for $w$, and therefore for $\varphi^j w$. Thus H82A1 proves that
\emph{every} such pair has finite class. All operations commute
with morphisms of the underlying lattice representations.

Suppose two cut vectors differ by $w'-a w\in vL_c$, with $a\in O$,
and take their two pairs. Subtract $a$ times the first pair from
the second. The truncation condition $\mathrm{IsTr}$ is additive
for this operation: take a common larger denominator $p^k$ for
the two witnesses and subtract them. For $u=w'-a w$ the resulting
witness $b\in vL_c$ satisfies
\[
 (\gamma-1)(b-p^ku)\in v^ML_c,\qquad
 p^k(\tau'-a\tau)-p^\sigma\varphi^j b\in D_j L(c,j).
\]
On $v^i L_c/v^{i+1}L_c$ the generator acts by
$\chi(\gamma)^{(p-1)i}$ for $i>0$, a scalar different from one.
There is no fixed vector on any such free grade, so induction
gives $b-p^ku\in v^M L_c$. The product formula for $\varphi^j(v)$
shows $(\varphi^j v)^M$ is divisible by $D_j$:
each $P_k$ occurs to power $(p-1)M\ge3g$ when
$M=\lfloor(h+2g)/(p-1)\rfloor+1$, $g\le h$.
It follows that $p^k(\tau'-a\tau-p^\sigma\varphi^ju)\in D_jL(c,j)$.
Cancel $p^k$ in this membership: $L(c,j)$ is free on the
Frobenius-translated basis, by invertibility after inverting its
cyclotomic factors. The nonzero reduction of $D_j$ modulo the
coefficient uniformizer makes its principal quotient torsion-free
over $O$ by divisibility in $O[[v]]$. This proves that the first
entry of the difference of the two pairs belongs to $D_jL(c,j)$.
For the $j$ thin letters after $c$ the Frobenius pole at each
step is at most $g$. Thus $D_jL(c,j)\subset v^ML_{c+j}$.

On the full cyclic lattice $v^M L$, Frobenius contracts:
$\varphi(v)$ is a unit times $vP_0^{p-1}$ and
$P_0^h\varphi L\subset L$, with $(p-1)M>h$.
Every iterate gains a factor in the Frobenius-stable ideal $(p,v)$,
so $\varphi-1$ is invertible on $v^M L$ by its geometric series.
Write the difference of the first entries as $(\varphi-1)d$
there. The cocycle law and injectivity imply the difference of
the second entries is $(\gamma-1)d$. The classes are therefore
equal up to the same scalar $a$. In particular the class is
independent of the truncation witness for a fixed cut vector.
% END H82A2

% BEGIN H82A3
\paragraph{H82A3. The integral congruence sufficient to lift each arc.}
Use the row B cuts and the actual onto map $f$ of I82B3. Denote
the physical cut spaces of H82A2 by $W_{1,c},W_{2,c}$ and their
lattices by $L_{1,c},L_{2,c}$. Write $R=rr_{\rm par}$,
$a_0=2\sigma+R+2g\ell_{\max}$ and $\kappa=\mathrm{kapPar}$.
The following sufficient test involves just one vector and reduction
modulo $v$ at a time. For $w\in W_{2,c}$ find
$w'\in W_{1,c}$, a unit $\varepsilon\in O^\times$ and an integer
$0\le k\le2\kappa$ such that
\[
 f(w')-\varepsilon p^{2\sigma+2a_0+R+k}w\in vL_{2,c}.
\]
It suffices to do this for the vectors $w_c$ produced from a target
extension in I82B5. Here is the precise implication to the native
row inequality.

For $w'$ take an integral truncation for $(\gamma-1)w'$ by
I82B4, let the first entry at $c+j$ be
$p^\sigma\varphi^j w'-\tau$, and solve for the second entry in
$v^M L_1$ using invertibility of $\varphi-1$ there. The difference
of the first entry belongs after application of $\gamma-1$ to
$v^M L_1$ by the last congruence of I82B4. These are therefore
an integral Herr pair of the required type. H82A1--H82A2 give
a finite class and show that its image is precisely
$\varepsilon p^{2\sigma+2a_0+R+k}$ times the target arc class.
Multiply by $\varepsilon^{-1}p^{2\kappa-k}$ and sum over cuts.
By I82B5 their unscaled target sum is $p^{R+2\sigma}$ times the
given finite target class. Thus the source finite image contains
the multiple at exponent
\[
    (2\sigma+2a_0+R+2\kappa)+(R+2\sigma)
       =2R+4\sigma+2a_0+2\kappa=e_{\rm par}.
\]
I82B1--I82B3 turn it into the required normalized length at this
row. In particular the target decomposition, the existence and
finiteness of the source arc classes, and their functorial comparison
consume the indicated exponents. The next chapters prove the displayed congruence in the actual
cut space of the given lattice map, culminating in B82R3.
% END H82A3
